% Options for packages loaded elsewhere
\PassOptionsToPackage{unicode}{hyperref}
\PassOptionsToPackage{hyphens}{url}
\PassOptionsToPackage{dvipsnames,svgnames,x11names}{xcolor}
\documentclass[
  11pt,
  a4paper,
]{article}
\usepackage{xcolor}
\usepackage[margin=24mm]{geometry}
\usepackage{amsmath,amssymb}
\setcounter{secnumdepth}{-\maxdimen} % remove section numbering
\usepackage{iftex}
\ifPDFTeX
  \usepackage[T1]{fontenc}
  \usepackage[utf8]{inputenc}
  \usepackage{textcomp} % provide euro and other symbols
\else % if luatex or xetex
  \usepackage{unicode-math} % this also loads fontspec
  \defaultfontfeatures{Scale=MatchLowercase}
  \defaultfontfeatures[\rmfamily]{Ligatures=TeX,Scale=1}
\fi
\usepackage{lmodern}
\ifPDFTeX\else
  % xetex/luatex font selection
\fi
% Use upquote if available, for straight quotes in verbatim environments
\IfFileExists{upquote.sty}{\usepackage{upquote}}{}
\IfFileExists{microtype.sty}{% use microtype if available
  \usepackage[]{microtype}
  \UseMicrotypeSet[protrusion]{basicmath} % disable protrusion for tt fonts
}{}
\makeatletter
\@ifundefined{KOMAClassName}{% if non-KOMA class
  \IfFileExists{parskip.sty}{%
    \usepackage{parskip}
  }{% else
    \setlength{\parindent}{0pt}
    \setlength{\parskip}{6pt plus 2pt minus 1pt}}
}{% if KOMA class
  \KOMAoptions{parskip=half}}
\makeatother
\usepackage{color}
\usepackage{fancyvrb}
\newcommand{\VerbBar}{|}
\newcommand{\VERB}{\Verb[commandchars=\\\{\}]}
\DefineVerbatimEnvironment{Highlighting}{Verbatim}{commandchars=\\\{\}}
% Add ',fontsize=\small' for more characters per line
\newenvironment{Shaded}{}{}
\newcommand{\AlertTok}[1]{\textcolor[rgb]{1.00,0.00,0.00}{\textbf{#1}}}
\newcommand{\AnnotationTok}[1]{\textcolor[rgb]{0.38,0.63,0.69}{\textbf{\textit{#1}}}}
\newcommand{\AttributeTok}[1]{\textcolor[rgb]{0.49,0.56,0.16}{#1}}
\newcommand{\BaseNTok}[1]{\textcolor[rgb]{0.25,0.63,0.44}{#1}}
\newcommand{\BuiltInTok}[1]{\textcolor[rgb]{0.00,0.50,0.00}{#1}}
\newcommand{\CharTok}[1]{\textcolor[rgb]{0.25,0.44,0.63}{#1}}
\newcommand{\CommentTok}[1]{\textcolor[rgb]{0.38,0.63,0.69}{\textit{#1}}}
\newcommand{\CommentVarTok}[1]{\textcolor[rgb]{0.38,0.63,0.69}{\textbf{\textit{#1}}}}
\newcommand{\ConstantTok}[1]{\textcolor[rgb]{0.53,0.00,0.00}{#1}}
\newcommand{\ControlFlowTok}[1]{\textcolor[rgb]{0.00,0.44,0.13}{\textbf{#1}}}
\newcommand{\DataTypeTok}[1]{\textcolor[rgb]{0.56,0.13,0.00}{#1}}
\newcommand{\DecValTok}[1]{\textcolor[rgb]{0.25,0.63,0.44}{#1}}
\newcommand{\DocumentationTok}[1]{\textcolor[rgb]{0.73,0.13,0.13}{\textit{#1}}}
\newcommand{\ErrorTok}[1]{\textcolor[rgb]{1.00,0.00,0.00}{\textbf{#1}}}
\newcommand{\ExtensionTok}[1]{#1}
\newcommand{\FloatTok}[1]{\textcolor[rgb]{0.25,0.63,0.44}{#1}}
\newcommand{\FunctionTok}[1]{\textcolor[rgb]{0.02,0.16,0.49}{#1}}
\newcommand{\ImportTok}[1]{\textcolor[rgb]{0.00,0.50,0.00}{\textbf{#1}}}
\newcommand{\InformationTok}[1]{\textcolor[rgb]{0.38,0.63,0.69}{\textbf{\textit{#1}}}}
\newcommand{\KeywordTok}[1]{\textcolor[rgb]{0.00,0.44,0.13}{\textbf{#1}}}
\newcommand{\NormalTok}[1]{#1}
\newcommand{\OperatorTok}[1]{\textcolor[rgb]{0.40,0.40,0.40}{#1}}
\newcommand{\OtherTok}[1]{\textcolor[rgb]{0.00,0.44,0.13}{#1}}
\newcommand{\PreprocessorTok}[1]{\textcolor[rgb]{0.74,0.48,0.00}{#1}}
\newcommand{\RegionMarkerTok}[1]{#1}
\newcommand{\SpecialCharTok}[1]{\textcolor[rgb]{0.25,0.44,0.63}{#1}}
\newcommand{\SpecialStringTok}[1]{\textcolor[rgb]{0.73,0.40,0.53}{#1}}
\newcommand{\StringTok}[1]{\textcolor[rgb]{0.25,0.44,0.63}{#1}}
\newcommand{\VariableTok}[1]{\textcolor[rgb]{0.10,0.09,0.49}{#1}}
\newcommand{\VerbatimStringTok}[1]{\textcolor[rgb]{0.25,0.44,0.63}{#1}}
\newcommand{\WarningTok}[1]{\textcolor[rgb]{0.38,0.63,0.69}{\textbf{\textit{#1}}}}
\usepackage{longtable,booktabs,array}
\newcounter{none} % for unnumbered tables
\usepackage{calc} % for calculating minipage widths
% Correct order of tables after \paragraph or \subparagraph
\usepackage{etoolbox}
\makeatletter
\patchcmd\longtable{\par}{\if@noskipsec\mbox{}\fi\par}{}{}
\makeatother
% Allow footnotes in longtable head/foot
\IfFileExists{footnotehyper.sty}{\usepackage{footnotehyper}}{\usepackage{footnote}}
\makesavenoteenv{longtable}
\setlength{\emergencystretch}{3em} % prevent overfull lines
\providecommand{\tightlist}{%
  \setlength{\itemsep}{0pt}\setlength{\parskip}{0pt}}
\usepackage{mathtools}
\allowdisplaybreaks
\setlength{\emergencystretch}{6em}
\DeclareUnicodeCharacter{2084}{\ensuremath{_4}}
\usepackage{bookmark}
\IfFileExists{xurl.sty}{\usepackage{xurl}}{} % add URL line breaks if available
\urlstyle{same}
\hypersetup{
  pdftitle={Formal verification of a universal bordered Jacobian identity over arbitrary commutative rings},
  colorlinks=true,
  linkcolor={blue},
  filecolor={Maroon},
  citecolor={Blue},
  urlcolor={blue},
  pdfcreator={LaTeX via pandoc}}

\title{Formal verification of a universal bordered Jacobian identity
over arbitrary commutative rings}
\author{}
\date{23 August 2026}

\begin{document}
\maketitle

\textbf{Formalization companion article, 23 August 2026}\\
\textbf{Status:} kernel-checked Lean 4 development; internally replayed;
not externally peer reviewed

\subsection{Abstract}\label{abstract}

Let \(A=\sum_{i=0}^r a_iX^i\) and \(B=\sum_{j=0}^s b_jX^j\) over a
commutative ring, and let \(M\) be the Jacobian of the coefficient
multiplication map \((A,B)\mapsto AB\). We give a Lean 4 formalization
of all signed maximal minors of \(M\) and of the determinant obtained by
appending an arbitrary row. The final theorems hold over every
commutative ring, including rings with zero divisors and the boundary
cases \(r=0\) or \(s=0\). The proof avoids division and rank
assumptions. It first proves an abstract cross-multiplied cofactor
relation over arbitrary rings, identifies one anchored minor with a
Sylvester resultant, cancels a single coefficient only in a universal
polynomial integral domain, and then specializes the universal identity
to the target ring. The pinned package builds without warnings or
unproved placeholders under Lean 4.32.1 and Mathlib v4.32.1. This is a
formal verification of the stated determinant identities, not a result
on the Jacobian or Hessian conjectures and not a priority claim for the
classical identity.

\subsection{1. Statement}\label{statement}

Fix \(r,s\ge0\) and a commutative ring \(R\). Put

\[
A(X)=\sum_{i=0}^r a_iX^i,
\qquad
B(X)=\sum_{j=0}^s b_jX^j.
\]

Order the coefficient variables as

\[
a_0,\ldots,a_r,b_0,\ldots,b_s.
\]

The coefficient multiplication map has \(r+s+1\) output coordinates and
\(r+s+2\) input coordinates. Its Jacobian is therefore a rectangular
matrix

\[
M\in\operatorname{Mat}_{r+s+1,r+s+2}(R),
\]

with entries

\[
M_{k,a_i}=b_{k-i},\qquad M_{k,b_j}=a_{k-j},
\tag{1.1}
\]

where out-of-range coefficients are zero. Define

\[
\kappa=(a_0,\ldots,a_r,-b_0,\ldots,-b_s)
\tag{1.2}
\]

and use the resultant convention

\[
\mathcal R_{r,s}(A,B)=\operatorname{Res}(B,A).
\tag{1.3}
\]

For a column \(k\), let \(M_{\widehat k}\) denote the order-preserving
submatrix obtained by deleting that column.

\subsubsection{Theorem 1.1 (signed maximal
minors)}\label{theorem-1.1-signed-maximal-minors}

For every commutative ring \(R\), all \(r,s\ge0\), every coefficient
family, and every column \(k\),

\[
\boxed{
(-1)^k\det M_{\widehat k}
=(-1)^{r(s+1)}\mathcal R_{r,s}(A,B)\kappa_k.}
\tag{1.4}
\]

\subsubsection{Theorem 1.2 (bordered
determinant)}\label{theorem-1.2-bordered-determinant}

Append an arbitrary row \(v=(v_0,\ldots,v_{r+s+1})\) below \(M\). Then

\[
\boxed{
\det\operatorname{border}(M,v)
=(-1)^{s(r+1)+1}\mathcal R_{r,s}(A,B)
\sum_kv_k\kappa_k.}
\tag{1.5}
\]

The Lean theorems are respectively
\texttt{BorderedJacobianUniversal.det\_mulJac\_submatrix} and
\texttt{BorderedJacobianUniversal.det\_bordered}.

\subsection{2. The ring-general cofactor
mechanism}\label{the-ring-general-cofactor-mechanism}

The identity

\[
M\kappa=0
\tag{2.1}
\]

is the differential form of \(AB-AB=0\). In coordinates, each output row
is the difference of the same convolution sum in opposite orders. The
formal proof first uses a sum-type index separating the \(a\)- and
\(b\)-columns, then transports the result along an explicit equivalence
to the manuscript order (1.1). This separation prevents an indexing
convention from being hidden inside simplification.

For a general \(n\times(n+1)\) matrix \(N\), define its signed cofactor
vector

\[
C_k=(-1)^k\det N_{\widehat k}.
\]

The familiar field argument would say that two kernel vectors are
proportional when the kernel is one-dimensional. That argument is
unusable over a ring with zero divisors: rank and division need not
behave as required. Instead the formalization proves the
cross-multiplied identity

\[
x_pC_k=x_kC_p
\tag{2.2}
\]

for every \(x\) with \(Nx=0\). The proof appends a selector row to
\(N\), uses the adjugate identity, and expands the resulting
determinants. It is valid over an arbitrary commutative ring and assumes
neither generic rank nor the invertibility of any coefficient.

The same abstract file proves Laplace expansion of an arbitrary border:

\[
\det\operatorname{border}(N,v)=(-1)^n\sum_kv_kC_k.
\tag{2.3}
\]

Thus Theorem 1.2 becomes a formal consequence of Theorem 1.1 once its
signs are reconciled.

\subsection{3. The anchored Sylvester
minor}\label{the-anchored-sylvester-minor}

Delete the column corresponding to \(a_r\). After separating its last
row, the remaining core is identified entry by entry with Mathlib's
Sylvester matrix. The formal theorem is

\[
\det M_{\widehat{a_r}}
=(-1)^{rs}\operatorname{Res}(B,A)a_r.
\tag{3.1}
\]

No positive-degree side condition is inserted: the proof includes
\(r=0\), \(s=0\), and \(r=s=0\). Combining (3.1) with the cofactor sign
gives

\[
C_{a_r}=(-1)^{r(s+1)}\operatorname{Res}(B,A)a_r.
\tag{3.2}
\]

The use of \(\operatorname{Res}(B,A)\), rather than
\(\operatorname{Res}(A,B)\), is part of the theorem statement and is
tested throughout the concordance. This avoids absorbing a
degree-dependent sign in informal notation.

\subsection{4. Universal cancellation and
specialization}\label{universal-cancellation-and-specialization}

Equation (2.2), with \(x=\kappa\) and \(p=a_r\), gives

\[
a_rC_k
=\kappa_k(-1)^{r(s+1)}\operatorname{Res}(B,A)a_r.
\tag{4.1}
\]

It is generally invalid to cancel \(a_r\) in \(R\). The proof therefore
moves to the universal ring

\[
U=\mathbb Z[A_0,\ldots,A_r,B_0,\ldots,B_s],
\]

an integral domain in which the universal variable \(A_r\) is nonzero.
The single cancellation in (4.1) is made only in \(U\), yielding the
universal cofactor identity (1.4).

For an arbitrary commutative ring \(R\), the coefficient assignment
defines a ring homomorphism \(U\to R\). The formal specialization layer
proves that this map commutes with every relevant matrix entry,
polynomial coefficient, resultant, determinant and cofactor. Mapping the
universal theorem into \(R\) therefore gives (1.4) without cancellation
in \(R\). Finally (2.3) gives (1.5).

This architecture is the central assurance gain of the formalization.
The arbitrary-ring theorem is not extrapolated from a generic-field
calculation; it is obtained by an explicit universal identity and
functorial specialization.

\subsection{5. Formal object and theorem
map}\label{formal-object-and-theorem-map}

{\def\LTcaptype{none} % do not increment counter
\begin{longtable}[]{@{}ll@{}}
\toprule\noalign{}
Mathematical object & Lean declaration \\
\midrule\noalign{}
\endhead
\bottomrule\noalign{}
\endlastfoot
coefficient extension by zero & \texttt{coeffExt} \\
multiplication Jacobian \(M\) & \texttt{mulJac} \\
kernel vector \(\kappa\) & \texttt{kappa} \\
abstract border & \texttt{border} \\
signed cofactor vector & \texttt{cofactorVec} \\
kernel identity (2.1) & \texttt{mulJac\_mulVec\_kappa} \\
cross relation (2.2) & \texttt{kernel\_cofactor\_proportional} \\
anchored minor (3.1) & \texttt{det\_anchorMinor} \\
universal cofactor formula & \texttt{universal\_cofactorVec} \\
arbitrary-ring maximal minors & \texttt{det\_mulJac\_submatrix} \\
arbitrary-ring border formula & \texttt{det\_bordered} \\
\end{longtable}
}

The detailed file-by-file and sign concordance is distributed as
\texttt{FORMALIZATION\_CONCORDANCE.md}.

\subsection{6. Boundary and zero-divisor
tests}\label{boundary-and-zero-divisor-tests}

Two executable formal tests address common failure modes.

\begin{enumerate}
\def\labelenumi{\arabic{enumi}.}
\tightlist
\item
  \texttt{det\_bordered\_zero\_zero} reduces the \(r=s=0\) theorem to
  the explicit \(2\times2\) determinant, checking the smallest dimension
  and its sign.
\item
  \texttt{zmod\_six\_one\_one\_control} instantiates a nonzero bordered
  determinant in \(\mathbb Z/6\mathbb Z\). It confirms that the released
  theorem is not accidentally restricted to domains or fields.
\end{enumerate}

Tests do not replace the general proof, but they expose off-by-one
indexing, sign and hidden-domain regressions at inexpensive concrete
points.

\subsection{7. Reproducibility receipt}\label{reproducibility-receipt}

The package pins:

\begin{itemize}
\tightlist
\item
  Lean 4.32.1, commit \texttt{f054605aea4b840552cca2e725580bffd1e1b704};
\item
  Lake 5.0.0;
\item
  Mathlib v4.32.1, commit
  \texttt{520045ab14e26149ee970e2e617ca04b09bde5d6}.
\end{itemize}

The accepted clean build processed 1,966 jobs with no warnings. A
fail-closed source scan found no \texttt{sorry}, \texttt{admit}, custom
\texttt{axiom}, \texttt{sorryAx}, \texttt{native\_decide},
\texttt{unsafe}, or \texttt{simp?}. The release-critical
\texttt{\#print\ axioms} output contains only the standard Lean/Mathlib
dependencies \texttt{propext}, \texttt{Classical.choice}, and
\texttt{Quot.sound}. The formal receipt has SHA-256

\begin{Shaded}
\begin{Highlighting}[]
\NormalTok{bc57d920fe4f6679a2b8fa871209f43de99de146c1c012b4a8a81dcadeb494a9}
\end{Highlighting}
\end{Shaded}

and binds every source, toolchain and concordance file by SHA-256. A
clean replay is

\begin{Shaded}
\begin{Highlighting}[]
\ExtensionTok{python3}\NormalTok{ verify\_release.py}
\end{Highlighting}
\end{Shaded}

from the formal package directory after fetching the pinned
dependencies.

\subsection{8. Trust and claim boundary}\label{trust-and-claim-boundary}

The successful build means that the stated Lean theorems elaborate and
are accepted by the pinned Lean kernel relative to Mathlib. It does not
independently verify the kernel, compiler, operating system or Mathlib.
The package has been internally replayed, not independently reproduced
or externally peer reviewed.

The formalization establishes only the determinant identities
(1.4)--(1.5). It does not prove invertibility of a polynomial map, the
two-dimensional Jacobian conjecture, the quartic Hessian conjecture, or
originality of the classical algebraic ingredients.

\subsection{Declarations}\label{declarations}

\textbf{Code availability.} Lean source, lockfiles, concordance,
verifier and the exact receipt accompany the successor package under the
Apache-2.0 license.

\textbf{Author responsibility.} The submitting author is responsible for
the mathematical statement, sign conventions and interpretation of the
formal result.

\textbf{AI-use disclosure.} AI systems contributed to proof
architecture, Lean implementation, replay, review and drafting. Internal
AI review is not external peer review.

\end{document}
